\begin{frame}{탐험에도 비용이 있다}
\begin{center}\begin{varwidth}{0.92\textwidth}\usebeamercolor[fg]{normal text}
  \textbf{무작위 팔의 기대 보상은 평균}
  \vskip0.5em
  \begin{center}
    탐험 스텝 하나당 기대 손실:
    \[q^*(a^*) - \frac{1}{k}\sum_a q^*(a)\]
  \end{center}
  
  \vskip0.4em
  \begin{itemize}
    
    
    \item 양쪽 끝이 나쁘면 \textbf{가운데 어딘가}에 최대치 존재.
  \end{itemize}
\end{varwidth}\end{center}
\note{\begin{itemize}
\item $\varepsilon$-greedy가 무작위 팔을 고르는 스텝에서, 그
  ``무작위 팔''의 기대 보상 = $k$개 팔의 평균
  $\frac{1}{k}\sum_a q^*(a)$ --- 최선의 팔 $q^*(a^*)$ 보다 \textbf{작다}.
\item $k=3$, $q^* = [1.0, 1.5, 2.0]$ 이면 스텝마다
  $2.0 - 1.5 = \mathbf{0.5}$ 의 기대 손실.
\item $\varepsilon$ \textbf{너무 작으면}: ``한 번의 나쁜 운''에
      최적 팔을 놓칠 확률 증가.
\item $\varepsilon$ \textbf{너무 크면}: 이미 찾은 최선 팔을
      외면한 채 스텝마다 0.5 손실 지속.
\end{itemize}}
\end{frame}
